% Generated draft: requires mathematical and source review.
\documentclass[11pt]{article}
\usepackage[margin=25mm]{geometry}
\usepackage{fontspec}
\setmainfont{Latin Modern Roman}
\usepackage{amsmath,amssymb,mathtools,unicode-math}
\setmathfont{Latin Modern Math}
\usepackage{enumitem,needspace,xurl,hyperref}
\setlength{\parindent}{0pt}
\setlength{\parskip}{5pt}
\setlength{\emergencystretch}{3em}
\newcommand{\N}{\mathbb{N}}
\newcommand{\Z}{\mathbb{Z}}
\newcommand{\Q}{\mathbb{Q}}
\newcommand{\R}{\mathbb{R}}
\newcommand{\C}{\mathbb{C}}
\newcommand{\F}{\mathbb{F}}
\newcommand{\A}{\mathbb{A}}
\newcommand{\PP}{\mathbb{P}}
\newcommand{\E}{\mathbb{E}}
\newcommand{\eps}{\varepsilon}
\newcommand{\dd}{\,\mathrm{d}}
\begin{document}
\section*{UnsolvedMath ID: 20001911}
\textbf{Source problem code:} catalogue entry
\subsection{Definitions and mathematical statement}
Let \((M^2,g)\) be a closed connected surface whose geodesic flow \(\phi^t\) on the unit tangent bundle \(SM\) is ergodic for normalized Liouville measure \(\mu_L\). The antecedents missing from this extract occur in companion record 20001910. For \(v\in SM\), lift its complete unit-speed geodesic to \(\widetilde\gamma_v:\R\to\widetilde M\), and define
\[
T_b=\{v:\widetilde\gamma_v(\R)\text{ is bounded}\},\quad
T_p=\{v:d(\widetilde\gamma_v(t),\widetilde\gamma_v(0))\to\infty
\text{ as }|t|\to\infty\},\quad
T_i=SM\setminus(T_b\cup T_p).
\]
The definitions are independent of the chosen lift. Thus \(T_p\) is the proper-lift set, while \(T_i\) consists of unbounded but nonproper lifts. These measurable invariant sets form a partition, so exactly one has full measure. The companion's extracted \(TM\) in the complement is interpreted as \(SM\).

Schmidt's Question 8.3.1 asks which set is conull. For orientable surfaces the stated expectation is \(\mu_L(T_i)=1\) on a torus and \(\mu_L(T_p)=1\) in genus at least two. For genus zero the universal cover is compact, so \(T_b=SM\). The passage suggests higher-dimensional analogues but does not specify a prediction there.

The supplied aid reports that, for every closed manifold with Liouville-ergodic geodesic flow, there is a deterministic \(\rho\in[0,1]\) such that for almost every \(v\),
\[
\lim_{t\to+\infty}\frac{d(\widetilde\gamma_v(t),\widetilde\gamma_v(0))}{t}
=\lim_{t\to-\infty}\frac{d(\widetilde\gamma_v(t),\widetilde\gamma_v(0))}{|t|}
=\rho.
\]
If \(\rho>0\), then \(\mu_L(T_p)=1\). For an ergodic Riemannian two-torus the aid reports vanishing homological Birkhoff drift by velocity reversal, hence \(\rho=0\), since deck displacement in \(\Z^2\) coarsely controls universal-cover distance. It also reports zero Liouville measure for globally minimizing vectors, whose lifts minimize distance between every pair of their points and have escape rate one.

Using a bounded measurable fundamental-domain section, let \(a_n(v)\in\Z^2\) be the deck-displacement cocycle at integer time. The remaining torus target is almost-sure unboundedness together with failure of proper escape, for example
\[
\sup_{n\in\Z}|a_n(v)|=\infty,\qquad
\liminf_{|n|\to\infty}|a_n(v)|<\infty.
\]
Zero drift alone proves neither recurrence nor unboundedness. The original minimizing-geodesic suggestion observes that an invariant minimizing set of positive measure would be conull by ergodicity, but also notes torus examples where that set has small Hausdorff dimension. The aid's zero-measure result shows that this positive-measure strategy cannot establish the torus prediction.

The appended section 8.4 defines the \(C^k\) topology on smooth metrics of a closed manifold by uniform control of coordinate components and derivatives through order \(k\), using finitely many compact coordinate cores covering the manifold. Different such choices give the same topology. This adjacent background supplies no additional generic-metric claim for Question 8.3.1.
\subsection{Short English statement}
For a surface with ergodic geodesic flow, are almost all lifted geodesics bounded, proper, or unbounded and recurrent? The source predicts the third behavior on tori and proper escape in higher genus. The supplied aid gives an escape-rate reduction, but zero torus drift alone does not prove the predicted recurrence and unboundedness.
\subsection{Sources}
\begin{itemize}
\item Schmidt, Question 8.3.1 in the AIM list \emph{Open Problems and Questions about Geodesics}: \url{https://aimath.org/pastworkshops/geodesicsproblems.pdf}.
\item Local registry \texttt{data/unsolvedmath/problems.json}, permanent ID 20001911, code \texttt{catalogue entry}; canonical \texttt{aim-geometry-notes.json}, \texttt{notes[248]}, source item 8.3; updated 14 August 2026. Companion permanent ID 20001910, code \texttt{AIM-GEOMETRY-0248}, \texttt{notes[247]}, supplies the preceding definitions and Liouville-ergodicity setting.
\item Supplied attempt 1: partial result, problem status at run open, rigor label proved, novelty confidence low; current status partially solved. The original literature field is empty. The aid is machine generated, not peer reviewed, and requires independent verification before citation; no detailed theorem references are supplied.
\end{itemize}
\end{document}
